% ===== §15 =====
\section{The Dialectical Sublation of Power}
\label{sec:powersublation}

\noindent
The synthesis has stated its definition, but it has not yet discharged the debt the thesis left it. The strongest of the traditions, the one that held power to be the very ontology of culture, raised an objection before the paper's central claim was even made: that there is no outside, no site free of power from which an endogenous culture could be generated, and that the claim of a private world between two people is therefore a naivety, since the intimacy in which the couple would build it is already shot through with the phallic organization of power (\xref{sec:power}). This objection must be met before the phenomenology of intimate culture can proceed, and it must be met, as promised, not by denial but by sublation---by conceding what is true in it and turning that concession into the ground of a different answer.

\subsection*{The concession}

The concession is total and unhedged: power \emph{is} internal to all culture. The paper does not claim that intimate relation is a sanctuary outside power, a pure space where two people generate a culture untouched by domination. This would be exactly the naivety the power tradition diagnosed, and the framework of the previous sections forbids it, for the same reason the tradition gives: culture spans the symbolic, the symbolic is organized by the phallic signifier, and the phallic signifier is the organ of power within the symbolic order. A culture that spans the symbolic therefore carries power in its very constitution. More than this: the process of generation itself produces power. Wherever value is generated, an asymmetry of contribution, of recognition, of dependence, is generated with it, and this asymmetry is power in its elementary form. There is no generation of culture, not even between two people who love each other, that does not generate power along with it. The psychoanalytic clear-sightedness the tradition shares---that there is no pure emancipatory space because there is no system of production and relation not already traversed by power---is accepted in full. Intimate relation generates power as surely as it generates culture, because generating the one is generating the other.

\subsection*{The turn: the objection was already answered by the fourth principle}

But the concession, which looks like a surrender, is the pivot of the reply, for the generation of power was already provided for in the ethics this series built before the objection was raised. The fourth of the five principles of just relational reproduction states that, since the generation of value carries power, a just generation must in the same movement generate the \emph{counter-power} that restrains power and prevents its entrenchment \citep{paperix}. The objection assumed that to concede power's internality was to concede the game---that if power is everywhere, then the dominated is everywhere and the endogenous culture is a fold of domination like any other. The fourth principle denies the inference. That power is internal to all generation does not entail that all generation is unjust; it entails that the justice of a generation lies not in the absence of power, which is impossible, but in whether the same process that generates power also generates the force that keeps that power from congealing. The question is not \emph{whether} there is power in the couple's culture---there always is---but whether the culture generates, alongside its power, the counter-power that prevents that power from entrenching into fixed domination. Justice is not purity; it is the internal generation of the restraint.

\subsection*{The mechanics: friction and the force that changes direction}

The mechanism can be given a form, and a mechanical image makes it exact. Power entering culture is like a body set moving on a surface. One cannot abolish the surface---one cannot generate culture without generating power, as one cannot have motion without a ground. But justice requires two forces to act upon the moving body. The first is a \emph{friction}, a counter-force that resists the unlimited accumulation of power and prevents it from running away into entrenched domination; without it, power accelerates along whatever gradient it finds and congeals into a fixed hierarchy, and the culture degrades into a mere reproduction of that hierarchy---the naturalized rank, the settled hegemony the power tradition described. The second is a \emph{force that changes direction}, a counter-force of non-entrenchment that keeps the structure re-orientable, able to be turned and re-formed rather than locked onto a single axis of authority. A culture with neither force is a body sliding down the gradient of its own power, and this is the fate the power tradition took for the destiny of all culture. A culture that generates both forces---friction against accumulation, and the power to change direction against fixation---is a culture in which power is present but not entrenched, exercised but not congealed. The two forces are not imported from outside the relation; they are generated within it, by the same process that generates the culture and its power, when that process is just.

\subsection*{The ontological deepening of power-to and power-over}

This gives the distinction between power-over and power-to, which an earlier paper of this series drew, an ontological depth it did not yet have \citep{paperxvii}. That distinction is often heard as a contrast between a bad power that dominates and a good power that enables, as though power-to were simply the absence of power-over, a benign capacity with no coercive edge. The present account shows that this is not quite right, and that the truth is stronger. Power-to is not the absence of power; it is power whose accumulation is met by an internally generated friction and whose direction is kept open by an internally generated force of re-orientation. Endogenous intimate culture is not a culture \emph{without} power---no culture is---but a culture in which power is \emph{non-dominating}, held in the condition of power-to by the continual internal generation of its own counter-power. The difference between a couple whose shared culture liberates and one whose shared culture dominates is not that the first has escaped power and the second has not; it is that the first generates, with its power, the friction and the re-orientation that keep the power from entrenching, and the second does not. With this, the objection of the power tradition is not evaded but answered on its own ground: yes, power is internal to all culture, including the most intimate; and precisely because the intimacy series built, into its ethics of relational reproduction, the requirement that just generation generate its own counter-power, the internality of power is not the end of the endogenous but the condition it must meet. The phenomenology of intimate culture can now proceed, no longer under the shadow of an objection it has not answered.
